\subsection{System \& Main IDE}
\label{chap:back_vscode}

The entire development environment used for this thesis utilizes Fedora 42\cite{FedoraProject2026} 
\cite{Fedora42Announcement2026}as the operating system. From the start, it was clear that a Linux-based operating 
system would be used for the thesis development, as Linux provides a native 
Unix-like toolchain that seamlessly supports STM32 firmware development and QEMU 
modifications. While Fedora offers certain advantages over other distributions by 
striking a balance between stability and incorporation of the latest open-source 
software, there is no specific reason to choose it over any other distribution beyond 
the developer's personal preference.

Unlike the operating system, the primary IDE used during development belongs 
to the Microsoft ecosystem. Visual Studio Code\cite{VisualStudioCode2024}, a lightweight, cross-platform, and 
open-source code editor, has been the main IDE used during the project's development 
due to its ease of use and widespread use in both professional and educational 
environments.